\section{Introduction}
\pagenumbering{arabic}
\subsection{Motivation}







In electronics there exists a trend towards ever smaller electrical components [2]. As the scale of electrical components shrinks, the quantum mechanical nature of the physics governing their behavior starts to dominate. There are many historical examples where quantum mechanics presents an obstacle for engineers when trying to miniaturize electrical components, but also many examples where quantum mechanics is leveraged to improve the function of small components or even produce new and hitherto impossible functionality [3]. Given this, it is easy to see the importance of understanding quantum mechanical charge transport. Understanding quantum mechanical charge transport requires doing experiments with systems in which quantum mechanical effects play a significant role. This is typically the case for systems that are either sufficiently small, or sufficiently cold [4]. The experiments detailed in this thesis are of this nature, examining the low temperature quantum mechanical phenomenon of superconductivity in a particular class of nanoscopic devices; electromigrated break junctions. \\


In certain materials transport becomes completely dissipationless at low enough temperatures [5]. Such materials are called superconductors, and the temperature at which the superconducting phase transition occurs is called the critical temperature, $T_C$. When a superconductor is cooled below its critical temperature, the electron-electron interaction becomes effectively attractive, which leads to the creation of paired electrons known as Cooper pairs. It is these Cooper pairs that are responsible for the superconducting transport. The Cooper pairs condense into a macroscopic ground state, and a gap in the density of states opens up around this ground state energy. This gap is referred to as the superconducting energy gap, $|\Delta|$.

When a superconductor and a normal metal are brought into contact, the superconductor can act as a reservoir injecting Cooper pairs into the normal metal. At temperatures close to the critical temperature of the superconductor, the Cooper pairs survive in the normal metal over a typical distance $\xi_N$ which is on the order of $\sqrt{D/T}$, where $D$ is the diffusion constant in the normal metal and $T$ is the temperature [6]. In this region, superconductivity is induced in the normal metal. This phenomenon is known as the superconducting proximity effect. At sufficiently low temperatures, $\xi_N$ becomes of the order of the size of the normal metal, and at sufficiently low energies all electrons in the normal metal will carry Cooper pair correlations. This means that, effectively, a gap is induced in the normal metal. The size of this gap is $\mathrm{min}\left(|\Delta|, \hbar D / L^2\right)$, where $L$ is the size of the normal layer. The quantity $\hbar D / L^2$ is called the Thouless energy [6].


\noindent This thesis specifically examines superconductivity and the superconducting proximity effect in superconducting electromigrated break-junctions (EMBJs), a specific class of break junction. Break junctions are devices consisting of two metal electrodes which are separated by a small gap. The gap which separates the electrodes in a break junction can be of the order of nanometers. Because of this, break junctions can be used as electrical contacts for nanoscopic objects, and their most prominent use is in electrically contacting molecules [8]. The ability of break junctions to facilitate single molecule molecular charge transport makes them particularly interesting, and motivates their study in this thesis. 

Superconductivity of the leads can add value to a single molecule transport measurement in many ways. Using superconducting leads can give otherwise hard to obtain spectroscopic information, to give one example. Superconducting leads can also proximitize the molecule and allow for it to be studied in the superconducting state. Since the charge in a superconductor is carried by Cooper pairs, while single molecule charge transport typically takes place one electron at a time, using superconducting leads allows one to study the competition between those effects [7]. These are some examples of the value that can be added by superconductivity.

Break junctions can be made by creating a gap in a wire, which can be done in multiple ways, including mechanically by stretching the wire until it breaks, chemically through etching, or in the case of EMBJs, by using electromigration. The different methods have different advantages and disadvantages. Among the different types of break junctions, EMBJs excel in terms of the achievable gate coupling [7]. When doing transport measurements    on molecules that behave as semiconductors, single electron transistors or quantum dots, for example, the presence of a gate electrode with sufficient coupling permits the tuning of channel doping and charge state, respectively. This is especially relevant in combination with superconductivity, since the high bias that is often required to reach the transport regime without gate tuning would cause superconducting electrodes to transition to the normal state. For these reasons, EMBJs are studied in this thesis.

EMBJs are made by creating a nanoscopic gap in a wire using electromigration. Electromigration is a process by which metal atoms are displaced that can take place at sufficiently high current densities. The dominant mechanism driving the displacement is the transfer of momentum between the electrons that make up the current and the atoms of the wire [1]. Electromigration is more likely to take place in the regions of the wire with a higher current density. By giving the wire a constricted geometry, such as the bowtie geometry used in this thesis, the constricted region will have a higher current density. This makes the location of the formed gap more consistent and controllable.  Electromigration is discussed in more detail in the chapter on theory.

Some materials are particularly well suited for the fabrication of EMBJs. Gold, for example, is a very popular material. Gold does not oxidize, which leads to cleaner interfaces. Gold also has a density of states that resembles that of a Fermi gas, which reduces the influence of the leads on the transport measurements. One drawback of gold is that it is not a superconductor. This drawback can be mitigated by inducing superconductivity in a gold based EMBJ using the superconducting proximity effect.

One particular type of molecule, the graphene nanoribbon (GNR), was specifically kept in mind during the design of the devices examined in this thesis. Narrow width GNRs have a large band gap and a notoriously high contact resistance when contacted using metal leads. This is due to a work function mismatch at the GNR contact interface [9], which can lead to Schottky barriers or poor band alignment. Making good electrical contact therefore requires choosing an appropriate metal to mitigate this work function mismatch as much as possible. If the interface resistance is not reduced, transport measurements will be dominated by the interface resistance. In that case, transparency will be low, which prevents the GNRs from becoming superconducting through the proximity effect. If the interface resistance is high enough, the current at low bias voltages and low temperatures becomes too small to be observed, and meaningful transport measurements are impossible. Palladium has been shown to be relatively well suited for contacting GNRs [10], which was kept in mind during the design of the devices studied in this thesis. Palladium is also not a superconductor, but like the example that was given above with gold, palladium too can be made superconducting using the proximity effect. \\

The goal of this thesis is to optimize the design and fabrication of proximity induced superconducting electromigrated break junctions. Optimize, in this context, means that the transport measurements can be done in the largest regime possible. In order to allow for superconducting experiments in the largest possible temperature range, $T_C$ should be as high as possible. In order to allow for superconducting charge transport in the largest possible bias range, $\mathrm{min}\left(E_{Th}, \Delta\right)$ should be as high as possible. The design should also be optimized such that electromigration is as consistent and controllable as possible. The design parameters that are studied are the material choice for the normal metal, the thicknesses of the normal metal and superconductor, and their geometry. The research question that this thesis tries to answer is the following. What are the effects of the bridge material and thickness, superconductor thickness, and the geometry of the device on the Thouless energy, critical temperature and quality of electromigration? 

During the experiments of this thesis, some of the devices showed an asymmetry in their critical current when a magnetic field was applied. In response to this finding, the research question "How does the asymmetry in critical current evolve as the junction is electromigrated to progressively higher resistance values?" was added.


\subsection{Structure of the Thesis}

Chapter \ref{section: Theory} of this thesis gives an overview of the relevant theory needed to understand the methods and results discussed in further chapters. Chapter \ref{section:Methods} contains a discussion of the methods used for the fabrication of the devices studied and the experiments done on them. In chapter \ref{section:Results} I report the results of the fabrication and experiments, followed by their discussion in chapter \ref{section:Discussion}. Chapter \ref{section:Conclusion} contains concluding remarks on the work of this thesis as well as a discussion on the outlook. 


